% Generated draft: requires mathematical and source review.
\documentclass[11pt]{article}
\usepackage[margin=25mm]{geometry}
\usepackage{fontspec}
\setmainfont{Latin Modern Roman}
\usepackage{amsmath,amssymb,mathtools,unicode-math}
\setmathfont{Latin Modern Math}
\usepackage{enumitem,needspace,xurl,hyperref}
\setlength{\parindent}{0pt}
\setlength{\parskip}{5pt}
\setlength{\emergencystretch}{3em}
\newcommand{\N}{\mathbb{N}}
\newcommand{\Z}{\mathbb{Z}}
\newcommand{\Q}{\mathbb{Q}}
\newcommand{\R}{\mathbb{R}}
\newcommand{\C}{\mathbb{C}}
\newcommand{\F}{\mathbb{F}}
\newcommand{\A}{\mathbb{A}}
\newcommand{\PP}{\mathbb{P}}
\newcommand{\E}{\mathbb{E}}
\newcommand{\eps}{\varepsilon}
\newcommand{\dd}{\,\mathrm{d}}
\begin{document}
\section*{UnsolvedMath ID: 700534}
\textbf{Source problem code:} catalogue entry
\subsection{Definitions and mathematical statement}
A group is locally nilpotent if every subgroup generated by finitely many elements is nilpotent. Here a group \(H\) is nilpotent if its lower central series, defined by \(\gamma_1(H)=H\) and \(\gamma_{n+1}(H)=[\gamma_n(H),H]\), reaches the trivial subgroup after finitely many steps; \([A,B]\) denotes the subgroup generated by the commutators \(a^{-1}b^{-1}ab\), where \(a\in A\) and \(b\in B\).

We use the standard Hausdorff convention for locally compact topological groups and interpret local nilpotence as the algebraic property just defined, without requiring finitely generated subgroups to be closed. The supplied record does not explicitly specify these conventions; they should be confirmed during mathematical review.

\textbf{Problem (V. S. Charin, 1982; described in the source as a well-known problem).} Let \(G\) be a locally compact, locally nilpotent topological group. Suppose that
\[
\text{for every subgroup }H\leq G,\qquad
\bigl(H\text{ is closed in }G\ \text{and}\ H\ne G\bigr)
\Longrightarrow H\text{ is compact in its subspace topology}.
\]
If \(G\) is non-compact, must it be abelian; that is, must
\[
xy=yx\qquad\text{for all }x,y\in G
\]
hold?
\subsection{Short English statement}
V. S. Charin asks whether every non-compact locally compact, locally nilpotent group whose proper closed subgroups are all compact must be abelian. The supplied registry labels the problem open; this draft does not independently verify that status.
\subsection{Sources}
\begin{itemize}
\item Local UnsolvedMath registry, supplied complete record: permanent ID \(700534\), exact source problem code \texttt{catalogue entry}, title \emph{Kourovka Notebook Problem 8.86}. The record supplies the status \texttt{open}.
\item \emph{The Kourovka Notebook (Unsolved Problems in Group Theory)}, Problem 8.86, attributed to V. S. Charin in the 8th issue (1982). The supplied background identifies the 21st issue edition (2026), \texttt{arXiv:1401.0300v48}; the collection description names E. I. Khukhro and V. D. Mazurov as editors. The statement is qualified as a well-known problem. This bibliographic attribution is supplied source information and has not been independently checked.
\end{itemize}
\end{document}
